\subsection{The global Tate module of the elliptic target}
\label{sec:ghost-elliptic-arithmetic}

Keep \(E':y^2=x^3+9\), \(E:y^2=x^3-243\),
the isogenies \(\phi:E'\to E\), \(\psi:E\to E'\),
and put
\[
                  P_0=(6,15)\in E'(\mathbb Q),
                  \qquad T=(0,3)\in E'(\mathbb Q).
\]

\begin{proposition}[The two exact three-isogeny Selmer calculations]
\label{prop:ghost-elliptic-selmer}
With the domain curve indicated by the isogeny, the
Selmer groups are
\[
 \operatorname{Sel}^{\phi}(E'/\mathbb Q)=0,
 \qquad
 \operatorname{Sel}^{\psi}(E/\mathbb Q)
      =\langle6,12\rangle\simeq(\mathbb Z/3)^2
\]
where the second group is a subgroup of
\(\mathbb Q^\times/(\mathbb Q^\times)^3\).
Both elliptic curves
have trivial three-primary Tate--Shafarevich groups,
and
\[
 \operatorname{rank} E'(\mathbb Q)=1,
 \qquad E'(\mathbb Q)^{\wedge3}
          =\mathbb Z_3 P_0\oplus(\mathbb Z/3)T.
\]
Here the completion is \(\varprojlim_n E'(\mathbb Q)/3^n\).
\end{proposition}
\begin{proof}
The kernels are \(E'[\phi]=\langle T\rangle\simeq
\mathbb Z/3\) and \(E[\psi]\simeq\mu_3\).
For the latter, the nonzero points are
\((0,\pm9\sqrt{-3})\), with cyclotomic Galois action.
The isogeny cohomology sequence and the local sequences
give the exact Selmer sequences
\[
 \begin{split}
 0&\longrightarrow E(\mathbb Q)/\phi E'(\mathbb Q)
 \longrightarrow\operatorname{Sel}^{\phi}(E'/\mathbb Q)
 \longrightarrow\operatorname{Sha}(E'/\mathbb Q)[\phi]
 \longrightarrow0,\\
 0&\longrightarrow E'(\mathbb Q)/\psi E(\mathbb Q)
 \longrightarrow\operatorname{Sel}^{\psi}(E/\mathbb Q)
 \longrightarrow\operatorname{Sha}(E/\mathbb Q)[\psi]
 \longrightarrow0.
 \end{split}                                             \tag{12}
\]
These follow by the same kernel diagram as the
multiplication sequence in
\cite[IV.2, equations (28)--(29)]{GhostMilneEC}:
replace multiplication by the stated isogeny, and
use its exact sequence of Galois modules. Thus the
source and target of each kernel in (12) are fixed.

Both curves have good reduction outside \(\{2,3\}\).
At such a prime the isogeny extends to a finite etale
isogeny of their smooth proper models, because its
degree three is invertible. A local point extends
over the valuation ring by properness; its inverse
image is therefore an unramified torsor. Consequently
every Selmer class is unramified outside \(2,3\).

For \(\phi\), the global candidate space is
\[
 H^1(G_{\mathbb Q,\{2,3,\infty\}},\mathbb Z/3)
                  =\langle\eta\bmod3\rangle.
\]
Indeed an order-three character is a cyclic cubic
abelian extension. Kronecker--Weber
\cite[I.4.16]{GhostMilneCFT} puts it in a cyclotomic
field of conductor supported on two and three.
The unit groups modulo powers of two have no
quotient of order three, and the three-primary
part of \((\mathbb Z/3^n)^\times\) is cyclic.
Its unique cubic character has conductor nine,
and is generated by \(\eta=(\chi^2-1)/24\bmod3\).

The local quotient
\(E(\mathbb Q_2)/\phi E'(\mathbb Q_2)\) is zero.
To see this without a reduction-type assumption,
the three-division polynomial of \(E\) is
\(3x(x^3-972)\). A point above \(x=0\) would
require a square root of \(-243\), which is
\(5\bmod8\). The other possibility would require
\(x^3=972\), impossible since \(v_2(972)=2\).
Thus \(E(\mathbb Q_2)[3]=0\).
This compact group has an open pro-two formal subgroup;
on that subgroup multiplication by three is invertible.
Its finite quotient cannot have three-torsion: an
order-three element there would lift uniquely to
an order-three element after correction in the
uniquely three-divisible subgroup. Hence multiplication
by three is an automorphism of \(E(\mathbb Q_2)\).
Since \(\phi\psi=[3]\), the isogeny \(\phi\)
is surjective on these local points. But
\(\eta(F_2)=(4-1)/24=1/8\ne0\bmod3\).
The only two nonzero global candidates thus fail
the local condition at two. This proves
\(\operatorname{Sel}^{\phi}=0\).

For \(\psi\), Kummer theory identifies the
candidate classes with cube classes supported on
two and three. This space has order nine, generated
by \(2,3\); the sign is a cube. The actual isogeny
Kummer map is
\[
          \delta_\psi(X,Y)=Y-3
                 \quad\bmod(\mathbb Q^\times)^3
                                                        \tag{13}
\]
away from its zero and pole. To verify the map rather
than merely its divisor, the formula for \(\psi\)
gives the literal identity
\[
 (Y-3)\circ\psi
            =\left(\frac{y-27}{3x}\right)^3.
\]
Conversely, a root \(r^3=Y-3\) determines
\(s=X/r\), with \(s^3=Y+3\), and
\[
 x=18/(s-r),\qquad y=27(s+r)/(s-r).
\]
These are the inverse equations of a \(\psi\)-fibre,
so (13) is its precise Kummer torsor. At the rational
points \(-T=(0,-3)\) and \(P_0=(6,15)\) its
classes are \(-6\), equal to \(6\) modulo cubes,
and \(12\). Their valuation vectors at \((2,3)\)
are \((1,1)\) and \((2,1)\), independent over
\(\mathbb F_3\). Thus rational points realize
all nine candidate classes. This proves the second
Selmer calculation, including every local condition,
and both right-hand kernels in (12) are zero.

The identities \(\psi\phi=[3]_{E'}\) and
\(\phi\psi=[3]_E\) now give
\(\operatorname{Sha}(E')[3]=\operatorname{Sha}(E)[3]=0\).
For example a class killed by three on \(E'\)
maps under \(\phi\) into the zero group
\(\operatorname{Sha}(E)[\psi]\), and hence belongs
to \(\operatorname{Sha}(E')[\phi]=0\).
Iterating multiplication by three kills every
three-primary class, so both full three-primary
subgroups are zero.

We also have \(E(\mathbb Q)=\phi E'(\mathbb Q)\),
and therefore
\(\psi E(\mathbb Q)=3E'(\mathbb Q)\).
Thus \(E'(\mathbb Q)/3E'(\mathbb Q)\) has order
nine, with basis the classes of \(-T,P_0\).
The three-division polynomial on \(E'\) is
\(3x(x^3+36)\), so its rational three-torsion
is exactly \(\langle T\rangle\). The finite-basis theorem
\cite[Chapter IV, printed p.~105]{GhostMilneEC}
then gives rank one. There is no rational
three-primary torsion of higher order, because it
would make \(T\) divisible by three, contradicting
its nonzero class in the displayed quotient.
The class of \(P_0\) generates the free part
modulo three, hence its three-adic completion,
by the elementary digit-lifting argument for a
finitely generated rank-one group. This proves
the last assertion.
\end{proof}

\begin{proposition}[Absolute continuous first cohomology of the target]
\label{prop:ghost-elliptic-h1}
The ordinary Kummer map is an isomorphism
\[
 E'(\mathbb Q)^{\wedge3}
       \xrightarrow{\ \sim\ }
 H^1_{\rm cts}(G_{\mathbb Q},T_3E')
       =\mathbb Z_3\kappa(P_0)
          \oplus(\mathbb Z/3)\kappa(T).
\]
This is absolute Galois cohomology, with no imposed
finite ramification set or local Selmer condition.
\end{proposition}
\begin{proof}
The finite Kummer sequences, with their coefficient
transition maps, give
\[
 0\longrightarrow E'(\mathbb Q)^{\wedge3}
 \longrightarrow H^1_{\rm cts}(G_{\mathbb Q},T_3E')
 \longrightarrow T_3H^1(\mathbb Q,E')
 \longrightarrow0.                                   \tag{14}
\]
The first finite tower has surjective transitions;
the comparison from continuous cochains has no
degree-one error because the finite invariant
groups have zero derived limit. On the last
groups \(H^1(\mathbb Q,E')[3^n]\), coefficient
reduction is multiplication by three, as follows
from the morphism of multiplication-by-\(3^n\)
Kummer sequences. We prove this actual last Tate
module zero.

Let \(\xi_n\in H^1(\mathbb Q,E')[3^n]\) satisfy
\(3\xi_{n+1}=\xi_n\). For each fixed finite
place \(q\ne3\), the three-primary part of
\(H^1(\mathbb Q_q,E')\) is finite. Local Tate
duality identifies it with the dual of
\(E'(\mathbb Q_q)^{\wedge3}\), which is finite
because the local group has an open pro-\(q\)
formal subgroup; see
\cite[I.3.3--3.5]{GhostMilneADT}.
A compatible three-division tower in a finite
three-group is zero: iterate the compatibility
past its exponent. Therefore every \(\xi_n\)
has zero localization at every \(q\ne3\).
The real localization also has zero three-primary
part. This argument applies to each place of the
entire compatible tower; it does not assume a
uniform ramification set for its individual terms.

The three-primary Tate--Shafarevich group was proved
zero in Proposition~\ref{prop:ghost-elliptic-selmer}.
We may therefore apply the three-primary exact
sequence of \cite[I.6.26(b)]{GhostMilneADT} with
all places. In particular, for each finite-order
global class \(\xi_n\), its local Tate pairings
with a rational point sum to zero. All terms
except the one at three are already zero, so
\[
                   \langle\xi_{n,3},P_0\rangle_3=0
                   \qquad\hbox{for every }n.          \tag{15}
\]
The point \(P_0\) is non-torsion globally by the
preceding rank-one calculation, and is non-torsion
in \(E'(\mathbb Q_3)\): a local equality
\([m]P_0=O\) would be the same rational-coordinate
equality over \(\mathbb Q\). That local compact
group has an open subgroup isomorphic to
\(\mathbb Z_3\); hence the closure of the cyclic
subgroup generated by \(P_0\) has finite index
in its pro-three completion. Its annihilator in
the three-primary local Weil--Chatelet group is
therefore finite, by the same local duality.
Equations (15) put the entire compatible local
tower in this finite group, so it too is zero.
Every \(\xi_n\) is now locally trivial at all
places, and hence is zero by the already proved
vanishing of the three-primary Tate--Shafarevich
group. Thus \(T_3H^1(\mathbb Q,E')=0\), and (14)
is the asserted isomorphism.
\end{proof}

\begin{proposition}[A support-prime test for every original A.44 output]
\label{prop:ghost-elliptic-support-test}
The image of full geometrically pro-three sections
of \(U\) whose abelian pairs are outputs of the
original A.44 recursion is contained in the following
explicit subset of those outputs:
\[
 \begin{split}
 &\operatorname{ord}_{E'(\mathbb F_{q_j})}((-2,1))
                              \text{ is prime to }3
                              \qquad\text{for every }q_j,\\
 &\operatorname{ord}_{E'(\mathbb F_{r_j})}((3,-6))
                              \text{ is prime to }3
                              \qquad\text{for every }r_j.
 \end{split}                                            \tag{16}
\]
Thus these tests apply to all choices of primes in
A.44, not only the prescribed examples of
Theorem~\ref{thm:ghost-original-first-integral-class3}.
\end{proposition}
\begin{proof}
Take an element of the stated section domain.
Proposition~\ref{prop:ghost-elliptic-target} proves
that its \(H_3\)-image is the point-three section
and it therefore factors through the pointed open
cover \(Y\). Push it through the compactification
\(Y\hookrightarrow E'\). The resulting proper
section is its continuous Tate-module class
\[
                 c=a\kappa(P_0)+b\kappa(T),
                   \quad a\in\mathbb Z_3,\quad
                   b\in\mathbb Z/3,
                                                        \tag{17}
\]
by Proposition~\ref{prop:ghost-elliptic-h1}.
Here the base section is the rational origin of
the pointed cover. Choosing a different lift to
that cover changes the origin sheet by a central
deck translation; formula (17) already permits
this torsion change, which disappears after
multiplication by three.

Let \(q\) be any new support prime. It is greater
than three. The local boundary-conjugacy theorem
makes the pushed-forward local section on \(U\)
the section of its ordinary point \(t_q\).
A lift of a point section through a finite etale
cover is a point section on that cover: its fixed
coset gives a rational sheet over the point, and
the corresponding decomposition subgroup gives
the unique lifted section up to the geometric
group of the cover. This is the same coset
comparison used in
Proposition~\ref{prop:ghost-elliptic-target}, now
applied to the full pro-three subgroup of index
twenty-seven. Thus the restriction of (17) is
\(\kappa_q(Z_q)\) for an actual
\(Z_q\in Y(\mathbb Q_q)\) above \(t_q\).

Write \(g_0:C\to E'\) for the elementary map
\((u,v)\mapsto(6uv,3(3u^3+2v^3))\). The exact
formulas (9)--(10) give
\[
 g_0=\psi\iota+P_0,
 \qquad
 g_0\iota^{-1}\phi=[3]_{E'}+P_0.                 \tag{18}
\]
These equalities are identities of morphisms of
smooth projective curves; the translations are
part of the formulas. All curves and isogenies
have good reduction at \(q>3\), and the cubic
cover extends there. If \(q=q_j\), the local
point reduces to \(t=0\), whose cubic points
have \(u=0\); their image under \(g_0\) is
\(-T\). If \(q=r_j\), it reduces to \(t=1\),
where \(v=0\) and that image is \(T\).
Consequently (18) gives the identities after
reduction
\[
 [3]\overline Z_q+\overline P_0=-\overline T
       \quad(q=q_j),\qquad
 [3]\overline Z_q+\overline P_0=\overline T
       \quad(q=r_j).
\]
The kernel of reduction is a pro-\(q\) formal
group and is uniquely three-divisible. Its
integral three-adic Kummer class is zero; hence
\[
 3c_q+\kappa_q(P_0)=-\kappa_q(T)\quad(q=q_j),
 \qquad
 3c_q+\kappa_q(P_0)=\kappa_q(T)\quad(q=r_j).
                                                        \tag{19}
\]
Substituting (17) gives
\((3a+1)\kappa_q(P_0)=\pm\kappa_q(T)\).
The scalar \(3a+1\) is a unit in \(\mathbb Z_3\).
Multiplication by three first shows
\(3\kappa_q(P_0)=0\); on this order-three
subgroup the scalar \(3a+1\) acts as one.
Thus (19) says exactly
\[
       \kappa_{q_j}(P_0+T)=0,
             \qquad\kappa_{r_j}(P_0-T)=0.             \tag{20}
\]
The group law gives
\(P_0+T=(-2,1)\) and \(P_0-T=(3,-6)\).
At a good prime distinct from three, the pro-three
completion of \(E'(\mathbb Q_q)\) is the
three-primary component of \(E'(\mathbb F_q)\),
by the same uniquely divisible reduction kernel;
compare \cite[IV.3.2]{GhostMilneEC}.
Therefore (20) is precisely the two finite point-order
conditions (16), proving the asserted containment
for the entire section domain.
\end{proof}

The universal exclusion problem now includes the
following explicit remaining arithmetic case:
the A.44 prime-choice outputs satisfying every
test (16). Proposition~\ref{prop:ghost-elliptic-passing}
constructs such outputs; no all-depth section has been
constructed for a member of this family. The calculation reduces the
constraint on a putative global proper section to
finite point orders at the actual support primes;
it does not infer a global equality from vanishing
at an arbitrarily sparse infinite set of primes.
